% !TeX root = ../../Book3.tex
% 纲要标题：### B.1 系数、表示与复杂度

\section{系数、表示与复杂度}
\label{b3:sec:B01}

计算理想或商环，首先需要把对象写成有限数据，并说明怎样检验输出是否代表原来的对象。本节以多项式、矩阵和分式为例，区分可直接执行的运算与仍须给出证明的结构判断。

% 纲要细目（与计划纲要同步）：
% * 多项式、局部化分式、幂级数截断和有限展示模的不同输入模型
% * 域上的 Gröbner 基与整环、主理想整环上的强 Gröbner 基；首系数的额外作用
% * 有效系数环所需的可判定性和运算条件；不声称任意 Noether 环都支持有效算法
% * 算术复杂度、次数增长、系数膨胀和最坏情形与常见实例的区别
%
% #### B.1.1 多项式、分式、幂级数截断与有限展示模的表示
% #### B.1.2 Gröbner 基、强 Gröbner 基与首系数
% #### B.1.3 有效系数环的可判定性与运算条件
% #### B.1.4 算术复杂度、次数增长与系数膨胀


\subsection{多项式、分式、幂级数截断与有限展示模的表示}
\label{b3:subsec:B01:01}

同一个代数对象，可以有适合证明的表示，也可以有适合计算的表示。例如，理想写成有限生成形式，并不等于已经能判定成员关系；有限维商环写出一个向量空间基，也还没有给出乘法。计算之前应先说明，输入包含哪些有限数据，允许执行哪些运算，答案又以什么形式返回。

\ProseParagraphBreak
以下以可进行精确运算的域 \(k\) 为系数域。多项式用非零系数与指数向量的有限表表示；稀疏表示保存非零项，稠密表示则保存给定次数范围内的全部系数。局部化 \(U^{-1}k[x]\) 的元素用一对多项式 \((f,u)\) 表示，其中 \(u\in U\)。这对数据还应说明分母属于哪个乘性集：在原点局部环中，条件是 \(u(0)\ne0\)；在 \(k[x]_h\) 中，分母可取 \(h\) 的幂。这两种分式不能仅凭同样的斜线记号混用。

\ProseParagraphBreak
任意形式幂级数通常没有有限的输入表。有限计算可以接收一个多项式截断、一个逐次提供系数的规则，或者满足某个方程并附带分支条件的特殊级数。前两者所能证明的结论不同：知道
\[
f\equiv f_N\pmod{(x_1,\ldots,x_n)^N}
\]
只确定次数小于 \(N\) 的系数，不能据此断言 \(f=0\)。例如零级数和 \(x_1^N\) 在这份数据中不可区分。

\begin{definition}\label{b3:def:B-presentation-data}
设 \(R\) 为具有指定有限数据表示的交换环，\(p,q\) 为非负整数。给定矩阵 \(A\in R^{q\times p}\)，本附录以
\[
R^p\xrightarrow{A}R^q\longrightarrow M\longrightarrow0
\]
作为模 \(M=\operatorname{coker}A\) 的展示数据；矩阵的列是关系，\(R^q\) 的标准基给出指定生成元。若 \(M\) 有分次，还须给出源、靶各基向量的次数，使每个矩阵项的次数与所表示的映射相符。
\end{definition}

给出 \(\operatorname{coker}A\) 无须断言左映射单射。两个展示描述同构模时，计算仍须记录生成元如何对应，否则从一个展示求出的关系不能直接当作另一个展示的关系。这一点在\secref{b3:sec:B04}中用矩阵等式具体说明。

\subsection{Gröbner 基、强 Gröbner 基与首系数}
\label{b3:subsec:B01:02}

域上的除法能把任意非零首系数变成 \(1\)。在整数上，这一步可能改变理想。设 \(R\) 为整环，并在 \(R[x_1,\ldots,x_n]\) 的单项式上固定全局项序。首单项式 \(\operatorname{LM}(f)\) 只记录幂；首项 \(\operatorname{LT}(f)\) 则包含系数。若
\[
\operatorname{LT}(g)=a x^\alpha,\qquad
\operatorname{LT}(f)=b x^\beta,
\]
用 \(g\) 一步消去 \(f\) 的首项，既要求 \(x^\alpha\mid x^\beta\)，又要求 \(a\mid b\)。

\begin{definition}\label{b3:def:B-strong-groebner}
设 \(R\) 为整环，\(I\subset R[x_1,\ldots,x_n]\) 为理想，项序为全局项序。若有限集 \(G\subset I\setminus\{0\}\) 满足：每个非零 \(f\in I\) 的首项都被某个 \(g\in G\) 的首项整除，则称 \(G\) 为 \(I\) 的\term{强 Gröbner 基}[strong Gröbner basis]。这里的整除同时包括系数整除和单项式整除。
\end{definition}

另一种较弱的条件是，所有 \(\operatorname{LT}(g)\) 生成 \(I\) 的首项理想。这两种条件在域上相同，在一般系数环上不同。例如在 \(\mathbb Z[x]\) 中，\(\{2x,3x\}\) 的首项生成 \((x)\)，但 \(x=3x-2x\) 的首项既不被 \(2x\) 整除，也不被 \(3x\) 整除。因此它不是 \((x)\) 的强 Gröbner 基，加入 \(x\) 后才得到单项首项整除的性质。

\ProseParagraphBreak
在有效主理想整环上，首系数的最大公因子及 Bézout 表示提供了额外的构造：若两个首项有系数 \(a,b\)，就可用 \(ua+vb=\gcd(a,b)\) 产生系数更小的首项。仅计算域上的 \(S\)-多项式会漏掉这一要求。本附录不建立整环上强 Gröbner 基的完整算法；后文的 FGLM、矩阵消元和签名说明均以域为系数环。

\begin{example}\label{b3:exm:B-integer-coefficients}
在 \(\mathbb Z[x]\) 中取 \(I=(2x,x^2)\)。集合 \(\{2x,x^2\}\) 是强 Gröbner 基：\(I\) 中次数为 \(1\) 的多项式，其 \(x\) 系数为偶数；次数至少为 \(2\) 的非零多项式，其首项被 \(x^2\) 整除。但 \(x\notin I\)。扩张到 \(\mathbb Q[x]\) 后却有 \(I\mathbb Q[x]=(x)\)。因此，先在分式域上求基再清除分母，不能自动恢复原整数理想的成员关系。
\end{example}

\subsection{有效系数环的可判定性与运算条件}
\label{b3:subsec:B01:03}

“系数环是 Noether 环”保证某些生成过程在抽象意义上有限，却没有告诉我们如何比较两个输入系数。有效算法首先需要有限可读的元素表示、相等判定、加法与乘法；在域上还需要非零元素的求逆。在一般系数环上，约化可能进一步需要判定一个系数是否属于给定的有限生成系数理想，并在属于时求出线性组合。

\ProseParagraphBreak
例如有理数可用互素整数对表示；整数的精确四则运算和最大公因子计算给出相等判定。有限域须指定其特征以及扩域时所用的不可约多项式。代数数域可以给出 \(\mathbb Q[T]/(p)\) 及 \(p\) 的不可约性依据，再把系数约化到次数小于 \(\deg p\)。若所关心的是实根中的某个嵌入，还需隔离区间等附加数据，但抽象域上的理想运算并不依赖选哪一个复嵌入。

\ProseParagraphBreak
浮点数的近似零判定不具备同样性质。一个很小的非零主元可以决定矩阵秩，也可以决定某个多项式的首项。精确 Gröbner 基证书需要确定这些系数是否为零；数值计算则须另有误差控制和认证方法。两类计算可以配合，但不能把“数值余项很小”作为理想成员的精确证明。

\begin{remark}
若输入系数是任意实数，而所谓表示仅允许不断给出近似值，则通常无法在有限时间内判定它是否恰为零。因此，“任意域上的定理”与“任意域上的可执行程序”是不同层次的陈述。下文说算法可执行时，始终包含了前述精确系数运算假设。
\end{remark}

\subsection{算术复杂度、次数增长与系数膨胀}
\label{b3:subsec:B01:04}

计算代价至少有三种来源：系数运算的次数、中间多项式出现的单项式数，以及每个系数所需的存储位数。以域运算计数，可以比较两种线性代数方法；若系数在 \(\mathbb Q\) 中，还必须考虑分子、分母的增长。同样次数和同样项数的多项式，系数为一位整数与为数千位整数时，实际计算量并不相近。

\begin{example}\label{b3:exm:B-bit-growth}
在 \(\mathbb Q[x_1,\ldots,x_r]\) 中考虑
\[
x_1-2,\quad x_2-x_1^2,\quad\ldots,\quad x_r-x_{r-1}^2.
\]
每个输入的次数至多为 \(2\)，系数很小，但唯一解为
\[
x_i=2^{\,2^{i-1}}.
\]
若把消元结果写为 \(x_r-2^{2^{r-1}}\)，常数项的二进制位数是 \(2^{r-1}+1\)。输出本身已要求指数增长的位数，因而不能仅凭输入次数小就期待小规模的显式答案。
\end{example}

次数增长也会扩大线性代数问题。\(n\) 个变量中次数恰为 \(d\) 的单项式有 \(\binom{n+d-1}{d}\) 个，次数至多为 \(d\) 的有 \(\binom{n+d}{d}\) 个。用一张系数矩阵同时处理许多约化时，这些数给出可能的列数。稀疏存储能避开最初不存在的系数，但行消元可能产生新的非零项。

\ProseParagraphBreak
因此，算法说明应把可证明的终止性、给定模型下的运算界和实际实例的表现分别报告。F4 型方法利用批量约化，FGLM 利用商代数已经有限维，模方法则控制有理系数的膨胀；它们改善的是不同部分，没有一种选择对所有输入都最好。
